\begin{tikzpicture}[
    % Definición de perspectiva isométrica limpia
    x={(0.866cm,0.5cm)}, y={(-0.866cm,0.5cm)}, z={(0cm,1cm)},
    scale=0.9,
    % ESTILOS UNIFICADOS CON LA PALETA UNADM
    link_robot/.style={draw=PPTGrey!40, line width=4pt, line cap=round},
    joint_robot/.style={ball color=white, opacity=0.9},
    % Estilo AABB (Gris/Naranja para indicar ineficiencia/error)
    aabb_style/.style={draw=PPTOrange, fill=PPTOrange!10, opacity=0.2, dashed, line width=0.6pt},
    % Estilo Bolas (Verde Esmeralda para indicar precisión/topología)
    sphere_vol/.style={ball color=PPTSteel, opacity=0.5},
    obstacle_style/.style={ball color=black!80},
    grid_style/.style={draw=PPTGrey!20, thin}
]

    % --- LADO IZQUIERDO: AABB (TRADICIONAL) ---
    \begin{scope}[xshift=0cm]
        \foreach \i in {-1,0,1,2} {
            \draw[grid_style] (\i,-1,0) -- (\i,2,0);
            \draw[grid_style] (-1,\i,0) -- (2,\i,0);
        }
        
        \node[font=\small\sffamily\bfseries, color=PPTGrey] at (0.5,0.5,-1) {AABB (Ingeniería Clásica)};
        
        % Robot
        \draw[link_robot] (0,0,0) -- (0,0,1.8) -- (1.5,0,2.8);
        
        % Articulaciones
        \fill[joint_robot] (0,0,0) circle (0.2);
        \fill[joint_robot] (0,0,1.8) circle (0.18);
        \fill[joint_robot] (1.5,0,2.8) circle (0.16);

        % Obstáculo flotante (punto negro)
        \fill[obstacle_style] (1.1,0.8,1.5) circle (0.12);

        % CAJA 1: Vertical
        \draw[aabb_style] (-0.4,-0.4,0) -- (0.4,-0.4,0) -- (0.4,0.4,0) -- (-0.4,0.4,0) -- cycle;
        \draw[aabb_style] (-0.4,-0.4,1.8) -- (0.4,-0.4,1.8) -- (0.4,0.4,1.8) -- (-0.4,0.4,1.8) -- cycle;
        \foreach \i/\j in {-0.4/-0.4, 0.4/-0.4, 0.4/0.4, -0.4/0.4} \draw[aabb_style] (\i,\j,0) -- (\i,\j,1.8);

        % CAJA 2: Inclinada
        \draw[aabb_style] (-0.3,-0.4,1.8) -- (1.6,-0.4,1.8) -- (1.6,0.5,1.8) -- (-0.3,0.5,1.8) -- cycle;
        \draw[aabb_style] (-0.3,-0.4,2.9) -- (1.6,-0.4,2.9) -- (1.6,0.5,2.9) -- (-0.3,0.5,2.9) -- cycle;
        \foreach \i/\j in {-0.3/-0.4, 1.6/-0.4, 1.6/0.5, -0.3/0.5} \draw[aabb_style] (\i,\j,1.8) -- (\i,\j,2.9);
            
        \node[PPTOrange, font=\footnotesize\sffamily\bfseries] at (0.5,1.8,3.2) {COLISIÓN FALSA};
    \end{scope}

    % --- LADO DERECHO: BOLAS ABIERTAS (TOPOLOGÍA) ---
    \begin{scope}[xshift=5.5cm]
        \foreach \i in {-1,0,1,2} {
            \draw[grid_style] (\i,-1,0) -- (\i,2,0);
            \draw[grid_style] (-1,\i,0) -- (2,\i,0);
        }

        \node[font=\small\sffamily\bfseries, color=PPTSteel] at (0.5,0.5,-1) {Recubrimiento (Topología)};

        % Esqueleto interno
        \draw[link_robot, opacity=0.3] (0,0,0) -- (0,0,1.8) -- (1.5,0,2.8);

        % Cadena de esferas (Más densas para evitar huecos)
        \foreach \z in {0.1, 0.4, 0.7, 1.0, 1.3, 1.6, 1.9}
            \fill[sphere_vol] (0,0,\z) circle (0.35);
        
        \foreach \t in {0.15, 0.3, 0.45, 0.6, 0.75, 0.9, 1.05}
            \fill[sphere_vol] ({1.5*\t}, 0, {1.8 + 1.0*\t}) circle (0.33);

        % Articulaciones
        \fill[joint_robot] (0,0,0) circle (0.2);
        \fill[joint_robot] (0,0,1.8) circle (0.18);
        \fill[joint_robot] (1.5,0,2.8) circle (0.16);

        % Obstáculo (Misma posición)
        \fill[obstacle_style] (1.1,0.8,1.5) circle (0.12);
        
        \node[PPTSteel!80!black, font=\footnotesize\sffamily\bfseries] at (0.5,1.8,3.2) {LIBRE (SIN COLISIÓN)};
    \end{scope}

\end{tikzpicture}